\documentclass[11pt]{article}
\usepackage[T1]{fontenc}
\usepackage{lmodern}
\usepackage{amsmath,amssymb,amsthm,mathtools}
\usepackage[hidelinks]{hyperref}
\usepackage{enumitem}
\newtheorem{theorem}{Theorem}[section]
\newtheorem{lemma}[theorem]{Lemma}
\newtheorem{corollary}[theorem]{Corollary}
\newtheorem{proposition}[theorem]{Proposition}
\newtheorem{definition}[theorem]{Definition}
\newtheorem{remark}[theorem]{Remark}
\newcommand{\gr}{\gamma_{\rm gr}}
\newcommand{\grt}{\gamma_{\rm gr}^{t}}
\newcommand{\timesdir}{\mathbin{\times}}

\title{Random counterexamples to three Grundy product conjectures}
\author{}
\date{August 2026}

\begin{document}
\maketitle

\begin{abstract}
We disprove the hypergraph expanding-sequence product conjecture and, by the
standard incidence-graph translations, the Grundy total-domination
direct-product and Grundy domination strong-product conjectures.  A
reflexive random tournament matrix of order $1937$ has, with positive
probability, no triangular witness of length $45$.  Its product with its
transpose nevertheless contains a $1937\times1937$ identity submatrix.
Translating this single probabilistic object gives a connected co-bipartite
graph counterexample and an incidence bipartite graph counterexample.  No
explicit order-$1937$ matrix is claimed, and we make no priority claim.
\end{abstract}

\section{Introduction}

The Grundy domination number is the maximum length of a legal sequence of
distinct vertices, where each term has a vertex in its closed neighborhood
not dominated by earlier terms.  The strong-product conjecture of
Bre{\v{s}}ar et al.~\cite{bresar2016} predicts
\(\gr(G\boxtimes H)=\gr(G)\gr(H)\); the analogous total-domination
conjecture predicts \(\grt(G\timesdir H)=\grt(G)\grt(H)\).
The hypergraph formulation and equivalence of these three conjectures were
given in~\cite{bresar2018} (published as~\cite{bresar2018doi}).  We provide a
counterexample to all three statements.  The construction repurposes the
classical random fooling-set/tournament and $A\otimes A^T$ mechanisms
\cite{pourmoradnasseri2016,friesen2014} for the expanding-sequence parameter.
Neither the random model nor the identity-block mechanism is new, and we make
no priority claim.

\section{Expanding sequences and matrices}

For a $0$--$1$ matrix $A$, write $A\otimes B$ for the Kronecker product.

\begin{definition}
For ordered injective row indices $r_1,\ldots,r_k$ and column indices
$c_1,\ldots,c_k$, call the pair a triangular witness if
\[
A[r_i,c_i]=1,\qquad A[r_i,c_j]=0\quad(i<j).
\]
Let $\tau(A)$ be the largest possible $k$.
\end{definition}

Viewing the rows as hyperedges on the column set, the condition says that
each new hyperedge contains a point not in the earlier hyperedges, with the
chosen private points ordered by the sequence.  Thus $\tau(A)$ is exactly the
maximum expanding-row (or expanding-hyperedge) sequence length.  In
graph-product terminology, the same matrix pattern is also called a
semi-induced matching; this should not be silently identified with every
definition of the standard semi-ladder index
\cite{chalermsook2013,fabianski2019}.

\begin{lemma}\label{lem:transpose}
$\tau(A^T)=\tau(A)$.
\end{lemma}
\begin{proof}
Transpose a triangular witness and reverse both orders.  The diagonal-one
conditions remain diagonal-one.  A zero $A[r_i,c_j]$ with $i<j$ becomes the
zero $A^T[c_j,r_i]$, which is above the diagonal after reversing the two
orders.  The operation is an involution.
\end{proof}

\section{The random theorem}

Let $M$ be the reflexive random tournament matrix of order $n$: its diagonal
is one, and for every unordered pair $\{u,v\}$ exactly one of $M[u,v]$ and
$M[v,u]$ is one, independently across pairs.

\begin{lemma}\label{lem:union}
For fixed injective ordered row and column $k$-tuples,
\[
\Pr[\text{they form a triangular witness}]\le 2^{-\binom{k}{2}}.
\]
\end{lemma}
\begin{proof}
There are $\binom{k}{2}$ upper-zero constraints.  If one is a diagonal cell,
the event is impossible because diagonal entries are one.  Two distinct
constraints cannot use the same directed cell, since both tuples are
injective.  If two constraints use the same unordered tournament edge in
opposite directions, they require both orientations to be zero, again making
the event impossible.  Therefore, whenever the event is nonempty, all
$\binom{k}{2}$ zero constraints use distinct independent tournament bits.
The diagonal-one constraints can only reduce the event.  Hence its probability
is at most $2^{-\binom{k}{2}}$.
\end{proof}

\begin{theorem}\label{thm:random}
For $n=1937$ there exists a reflexive tournament matrix $M$ with
$\tau(M)\le44$ and
\[\tau(M\otimes M^T)\ge1937.
\]
\end{theorem}
\begin{proof}
Union bound over the $(n)_k^2$ ordered injective row and column tuples gives
\[
 \Pr[\tau(M)\ge k]\le (n)_k^2 2^{-\binom{k}{2}}
 < n^{2k}2^{-\binom{k}{2}}.
\]
For $n=1937$, $k=45$, $n<2^{11}$ and $\binom{45}{2}=990$, so the last
quantity is less than
\[
(2^{11})^{90}2^{-990}=2^0=1.
\]
Thus some $M$ has $t:=\tau(M)\le44$.  For rows $(i,i)$ and columns $(j,j)$,
\[
(M\otimes M^T)[(i,i),(j,j)]=M[i,j]M[j,i]=\delta_{ij}.
\]
These rows and columns form an identity submatrix, whose diagonal witness has
length $n=1937$.  This proves the second assertion.
\end{proof}

Theorem~\ref{thm:random} is already a counterexample to multiplicativity of
the hypergraph expanding-sequence parameter, since
\[
\tau(M\otimes M^T)\ge1937>44^2\ge\tau(M)\tau(M^T).
\]

\section{Strong-product graph corollary}

For a graph $G$, a closed-neighborhood sequence is legal if each selected
vertex has a private vertex in its closed neighborhood outside the union of
the earlier closed neighborhoods.  The Grundy domination number $\gr(G)$ is
the maximum length of such a sequence.  We use
$N_{G\boxtimes H}[(g,h)]=N_G[g]\times N_H[h]$.

Define $K(M)$ as the co-bipartite graph with cliques
$R=\{r_i:i\in[n]\}$ and $C=\{c_j:j\in[n]\}$, with cross-edge $r_ic_j$ iff
$M[i,j]=1$.

\begin{lemma}\label{lem:K}
\(\gr(K(M))=\tau(M)\).
\end{lemma}
\begin{proof}
An $R$-only legal sequence is exactly an expanding-row sequence of $M$;
similarly a $C$-only sequence is one of $M^T$.  By Lemma~\ref{lem:transpose}
both have maximum length $\tau(M)$.  If a legal sequence switches cliques,
then after the first selected vertex from each clique both cliques are
dominated, so the switch is its last term.  Suppose the prefix lies in $R$
and has length $s$, and let $c_j$ be the switching term.  Its private witness
is either in $C$ or in $R$.  If it is a column $c_x$, then $c_x$ was not
covered by the prefix.  Hence $r_x$ was not selected (otherwise reflexivity
would cover $c_x$), and $r_x$ can be appended using $c_x$ as private witness.
If it is a row $r_x$, then $r_x$ was not selected and its reflexive closed
neighborhood supplies the same witness when $r_x$ is appended.  Thus every
mixed sequence can be replaced by an $R$-only legal sequence of at least the
same length.  The transpose argument handles a prefix in $C$.  Conversely,
an inclusion-maximal expanding sequence is dominating: an undominated column
$c_j$ cannot have its index among the selected rows, so the unselected row
$r_j$ is appendable; the row case is identical.  Hence the maximum is
$\tau(M)$.
\end{proof}

The graph $K(M)$ is connected, co-bipartite, and has $3874$ vertices.  In
$K(M)\boxtimes K(M)$ put $z_i=(r_i,c_i)$ and $w_j=(c_j,r_j)$.  Since
$N[(u,v)]=N[u]\times N[v]$,
\[
w_j\in N[z_i]\quad\Longleftrightarrow\quad M[i,j]M[j,i]=\delta_{ij}.
\]
Thus $z_1,\ldots,z_n$ is legal, with private witnesses $w_1,\ldots,w_n$.
Any finite legal prefix extends to a dominating sequence by appending legal
vertices, so
\[
\gr(K(M)\boxtimes K(M))\ge1937>44^2\ge\gr(K(M))^2.
\]

\section{Total-Grundy direct-product corollary}

For a graph $G$ without isolated vertices, an open-neighborhood sequence is
legal if each term has a vertex in its open neighborhood outside the earlier
open neighborhoods; the maximum length is the Grundy total-domination number
$\grt(G)$.  The direct product has vertex set $V(G)\times V(H)$ and adjacency
$(g,h)\sim(g',h')$ exactly when $gg'\in E(G)$ and $hh'\in E(H)$.  The
hypergraph parameter above is the corresponding expanding-edge parameter.

For a matrix $A$, let $B(A)$ be its incidence bipartite graph: left vertices
are rows, right vertices are columns, and $i\sim j$ iff $A[i,j]=1$.

\begin{lemma}\label{lem:incidence}
\[
\grt(B(A))=\tau(A)+\tau(A^T)=2\tau(A).
\]
\end{lemma}
\begin{proof}
The open neighborhoods of left-side terms lie in the right side, while those
of right-side terms lie in the left side, so the two subsequences do not
interfere.  Each side is exactly a triangular witness, giving the first
equality; Lemma~\ref{lem:transpose} gives the second.
\end{proof}

The four vertex types in $B(M)\timesdir B(M)$ split into two disjoint
subgraphs by whether the two coordinates lie on the same side.  Directly from
the incidence condition, with the displayed coordinate orderings,
\[
B(M)\timesdir B(M)\cong B(M\otimes M)\;\sqcup\;B(M\otimes M^T).
\]
The first isomorphism uses $(R,R)$ as rows and $(C,C)$ as columns; the second
uses $(R,C)$ as rows and $(C,R)$ as columns.  Grundy total-domination is
additive over disjoint unions, because open neighborhoods do not cross the
union.
The second factor is $M\otimes M^T$, not another copy of $M\otimes M$.
The standard lexicographic witness gives
$\tau(M\otimes M)\ge t^2$: take the witness pairs for $M$ in outer order and
use the witness pairs in inner order within each outer pair.  The identity
submatrix in Theorem~\ref{thm:random} gives
$\tau(M\otimes M^T)\ge n$.  Therefore, using disjoint components additively,
\[
\grt(B(M)\timesdir B(M))
\ge2t^2+2n>4t^2=\grt(B(M))^2.
\]

\section{Asymptotic strengthening}

Fix $\varepsilon>0$ and set
$k=\lceil(4+\varepsilon)\log_2 n\rceil$.  The base-two logarithm of the
union bound is
\[
2k\log_2n-\frac{k(k-1)}2
=-(2\varepsilon+\varepsilon^2/2)(\log_2n)^2+O(\log n),
\]
which tends to $-\infty$.  Hence for all sufficiently large $n$ some
tournament matrix satisfies $\tau(M)=O(\log n)$, while
$\tau(M\otimes M^T)\ge n$.  The multiplicative violation is therefore
$\Omega(n/\log^2n)$.  The two graph translations above inherit the same
unbounded gap (up to the fixed square/product conversion).

\section{Related work and scope}

The strong-product conjecture originated in~\cite{bresar2016}.  The
total-Grundy direct-product conjecture, hypergraph formulation, and stated
equivalence appear in~\cite{bresar2018,bresar2018doi}.  The forest-factor
positive theorem is~\cite{bell2021}.  A proof attempt was later withdrawn;
the arXiv record explicitly marks version 2 as withdrawn and notes an error
in its proof~\cite{herrman2023}.

The random fooling-set/tournament perspective and the classical tensor
identity mechanism belong to the fooling-set/rank literature
\cite{pourmoradnasseri2016,friesen2014}; the expanding-sequence/semi-induced-
matching correspondence also appears in graph-product work
\cite{chalermsook2013}.  Standard semi-ladder terminology is surveyed in
\cite{fabianski2019} and recent complexity work~\cite{manurangsi2026}.
Our computational package exhaustively checks only small tournaments for
mutation and indexing errors; it is not evidence for the existential
$n=1937$ theorem.  The literature search recorded in
the file `research/LiteratureAudit.md` does not establish priority for this
counterexample.

\bibliographystyle{plain}
\bibliography{references}
\end{document}
